\section{Execution Diagnostics and Interval Sensitivity}
\label{app:review_diagnostics}

Tables~\ref{tab:grid_modes} and~\ref{tab:grid_geometry} use the same
locked-mode, final-policy returns as the main comparison and the Ant
development study. The utility and shortfall means use equal seed
weights; modes and requests within each reported stratum are equally
weighted within a seed. The same-reset command maximum is descriptive.

% Generated by make_evidence_tables.py from completed readouts.
\begin{table}[!htbp]
\centering\small\setlength{\tabcolsep}{4.5pt}
\caption{Interval sensitivity for the same original paired differences. Both intervals are marginal 97.5\%; the seeded percentile bootstrap was specified in the original protocols; the paired $t$ interval is a sensitivity analysis.}
\label{tab:ci_sensitivity}
\begin{tabular}{@{}llrrr@{}}
\toprule
Task & Utility contrast & Mean & Percentile bootstrap & Paired $t$ \\
\midrule
HalfCheetah & BAMOR -- native & $+105.15$ & ${[-33.80,235.83]}$ & ${[-44.20,254.50]}$ \\
HalfCheetah & BAMOR -- full-simplex & $+14.13$ & ${[-99.42,121.50]}$ & ${[-108.43,136.70]}$ \\
Hopper & BAMOR -- native & $+24.67$ & ${[-58.64,105.58]}$ & ${[-67.79,117.13]}$ \\
Hopper & BAMOR -- full-simplex & $+83.17$ & ${[3.03,165.15]}$ & ${[-6.80,173.14]}$ \\
\bottomrule
\end{tabular}
\end{table}

% Generated by make_evidence_tables.py from completed readouts.
\begin{table}[!htbp]
\centering\small\setlength{\tabcolsep}{4.5pt}
\caption{Requested utility $U$ and same-actor command shortfall $G$ by locked mode. Means over 20 training seeds on HalfCheetah and Hopper and five on Ant.}
\label{tab:grid_modes}
\begin{tabular}{@{}llrrr@{}}
\toprule
Task & Locked mode & BAMOR $U/G$ & Native $U/G$ & Full-simplex $U/G$ \\
\midrule
HalfCheetah & Normal & $357.8/208.4$ & $231.2/379.7$ & $416.2/338.8$ \\
HalfCheetah & Low speed & $-17.5/40.2$ & $-261.2/232.8$ & $-188.0/222.8$ \\
HalfCheetah & Low cost & $720.4/385.9$ & $612.5/533.0$ & $629.2/565.5$ \\
HalfCheetah & Locked & $124.6/91.0$ & $182.1/70.7$ & $271.3/67.8$ \\
Hopper & Normal & $523.1/347.1$ & $519.3/419.4$ & $452.1/442.0$ \\
Hopper & Fragile & $483.8/265.8$ & $463.1/315.5$ & $392.1/376.5$ \\
Hopper & Cheap energy & $362.4/298.3$ & $333.0/339.6$ & $290.1/330.9$ \\
Hopper & Grounded & $377.0/153.1$ & $332.3/234.6$ & $279.4/268.5$ \\
Ant & Energy save & $-329.6/1082.4$ & $696.8/737.6$ & -- \\
Ant & Race & $1670.6/40.7$ & $1681.4/58.7$ & -- \\
Ant & Cautious & $1625.5/59.6$ & $1716.5/51.4$ & -- \\
\bottomrule
\end{tabular}
\end{table}

% Generated by make_evidence_tables.py from completed readouts.
\begin{table}[!htbp]
\centering\small\setlength{\tabcolsep}{4.5pt}
\caption{Requested utility $U$ and same-actor shortfall $G$ by region of the requested preference. The maximum always searches the full grid; the region restricts only requests. Boundary excludes vertices.}
\label{tab:grid_geometry}
\begin{tabular}{@{}llrrrr@{}}
\toprule
Task & Region & Weights & BAMOR $U/G$ & Native $U/G$ & Full-simplex $U/G$ \\
\midrule
HalfCheetah & Vertex & 2 & $516.8/59.3$ & $162.4/434.1$ & $378.1/330.1$ \\
HalfCheetah & Interior & 19 & $273.1/194.2$ & $194.2/290.4$ & $272.1/295.4$ \\
Hopper & Vertex & 3 & $451.2/268.0$ & $316.2/446.8$ & $258.2/469.1$ \\
Hopper & Boundary & 12 & $433.9/269.7$ & $389.6/351.0$ & $321.4/387.7$ \\
Hopper & Interior & 6 & $434.6/257.8$ & $504.4/220.0$ & $465.1/230.7$ \\
Ant & Vertex & 3 & $1314.7/71.0$ & $1074.2/578.3$ & -- \\
Ant & Boundary & 12 & $1004.0/379.5$ & $1360.9/287.2$ & -- \\
Ant & Interior & 6 & $795.6/585.4$ & $1518.1/125.5$ & -- \\
\bottomrule
\end{tabular}
\end{table}

% Generated by make_evidence_tables.py from completed readouts.
\begin{table}[!htbp]
\centering\small\setlength{\tabcolsep}{4.5pt}
\caption{Ant mode 0 among requests with $0 < w_{\mathrm{energy}} < 1$ (five seeds). $\Delta U$ is BAMOR minus native CAPQL. Energy wins compare the same BAMOR actor queried at the pure-energy vertex versus the requested weight, both scored at the request.}
\label{tab:ant_geometry}
\begin{tabular}{@{}lrrr@{}}
\toprule
Region & Seed--weight cells & $\Delta U$ & Energy wins \\
\midrule
Boundary & 40 & $-1495.1$ & 37/40 \\
Interior & 30 & $-1982.9$ & 29/30 \\
\bottomrule
\end{tabular}
\end{table}

\FloatBarrier

\subsection{Allocation traces for the 75k score ablation}

Count displacement gives the minimum number of post-warmup assignments
that must change to make the two slot histograms equal. The chronological
measure counts mismatches at equal block indices, irrespective of whether
the two runs selected the same slots in another order. Slot names follow
Table~\ref{tab:operators}.
% Generated by make_evidence_tables.py from completed readouts.
\begin{table}[!htbp]
\centering\small\setlength{\tabcolsep}{4.5pt}
\caption{Allocation change within seed after the first ten deterministic blocks. Count displacement is half of the L1 distance between the two complete-block count vectors; order changes compare slots at equal block index. Means across 20 pairs.}
\label{tab:slot_shift}
\begin{tabular}{@{}lrrr@{}}
\toprule
Task & Count displacement & Order changes & Partial slot differs \\
\midrule
HalfCheetah & $6.2/36$ & $12.9/36$ & $12/20$ \\
Hopper & $8.0/36$ & $17.3/36$ & $15/20$ \\
\bottomrule
\end{tabular}
\end{table}

% Generated by make_evidence_tables.py from completed readouts.
\begin{table}[!htbp]
\centering\footnotesize
\caption{Mean complete feedback blocks per allocation slot at 75k across 20 seeds per task. Each cell is score-on / score-neutral. Post-warmup counts exclude the ten warmup blocks. Neither set counts the final 1,000-step block. HC denotes HalfCheetah.}
\label{tab:slot_counts}
\begin{tabular}{@{}lrrrr@{}}
\toprule
Slot & HC all & HC post-warmup & Hopper all & Hopper post-warmup \\
\midrule
C1 & $2.70/2.80$ & $1.70/1.80$ & $2.75/3.20$ & $1.75/2.20$ \\
C2 & $3.15/3.15$ & $2.15/2.15$ & $2.90/2.55$ & $1.90/1.55$ \\
C3 & $2.00/2.10$ & $2.00/2.10$ & $2.45/1.85$ & $2.45/1.85$ \\
C4 & $2.05/2.10$ & $2.05/2.10$ & $1.80/2.20$ & $1.80/2.20$ \\
C5 & $2.50/2.00$ & $2.50/2.00$ & $1.85/1.95$ & $1.85/1.95$ \\
C6 & $2.90/3.00$ & $1.90/2.00$ & $3.20/3.20$ & $2.20/2.20$ \\
C7 & $2.75/2.55$ & $1.75/1.55$ & $3.25/3.00$ & $2.25/2.00$ \\
R1 & $2.40/2.65$ & $1.40/1.65$ & $2.55/2.15$ & $1.55/1.15$ \\
R2 & $1.60/1.75$ & $1.60/1.75$ & $1.35/1.40$ & $1.35/1.40$ \\
R3 & $1.80/1.35$ & $1.80/1.35$ & $1.55/1.70$ & $1.55/1.70$ \\
R4 & $1.45/1.65$ & $1.45/1.65$ & $1.80/2.10$ & $1.80/2.10$ \\
R5 & $2.55/2.75$ & $1.55/1.75$ & $2.65/2.95$ & $1.65/1.95$ \\
R6 & $1.55/1.40$ & $1.55/1.40$ & $1.50/1.75$ & $1.50/1.75$ \\
R7 & $2.30/2.45$ & $1.30/1.45$ & $2.25/2.15$ & $1.25/1.15$ \\
E1 & $3.20/3.10$ & $2.20/2.10$ & $4.00/3.95$ & $3.00/2.95$ \\
E2 & $3.75/4.00$ & $2.75/3.00$ & $3.40/3.70$ & $2.40/2.70$ \\
E3 & $4.30/4.00$ & $3.30/3.00$ & $3.70/3.35$ & $2.70/2.35$ \\
E4 & $3.05/3.20$ & $3.05/3.20$ & $3.05/2.85$ & $3.05/2.85$ \\
\bottomrule
\end{tabular}
\end{table}

\FloatBarrier

\section{Secondary Archive Coverage and 50k Routing Study}
\label{sec:archive_quota}

In the separate 50k cohorts, archive hypervolume exceeded final-actor
hypervolume on all 25 seeds, with ratios of means of $+35.10\%$ on
HalfCheetah, $+91.53\%$ on Hopper and $+22.91\%$ on Ant
(Table~\ref{tab:matched_archive}). The final actor is one of the 19
snapshots, evaluated with the same final encoder. Weak monotonicity follows
from inclusion (Proposition~\ref{prop:hv_current}); the measured increases
reflect additional covered behaviors, not an evaluated deployment rule.

% Generated by make_evidence_tables.py from completed readouts.
\begin{table}[t]
\centering\small
\caption{Retrospective 50k comparison: v25 minus CAPQL, with 19 actors in each archive. Percentages compare cohort means; W/L/T denotes paired seed wins/losses/ties. The early CAPQL runs omitted two RNG seeds (Section~\ref{sec:limits}).}
\label{tab:matched_archive}
\begin{tabular}{lrrrrr}
\toprule
Environment & $n$ & Archive HV & W/L/T & Final HV & W/L/T \\
\midrule
HalfCheetah & 15 & +16.21\% & 10/5/0 & +0.74\% & 8/7/0 \\
Hopper & 5 & +140.39\% & 5/0/0 & +58.67\% & 3/2/0 \\
Ant & 5 & +1.17\% & 3/2/0 & -11.07\% & 0/5/0 \\
\bottomrule
\end{tabular}
\end{table}


The 50k score-neutral comparison retained quotas, warmup and exploration.
Mean archive and final hypervolume differences were negative on HalfCheetah
and Ant. Hopper's mean Final HV rose by $59.56\%$ with scores (two wins,
one loss, two ties), but omitting the largest favorable pair gives
$-2.96\%$. Table~\ref{tab:selector_effect} retains all five pairs.

% Generated by make_evidence_tables.py from completed readouts.
\begin{table}[t]
\centering\small
\caption{Feedback attribution at 50k: unchanged v25 adaptive minus its score-blind control on the same cohorts as Table~\ref{tab:matched_archive}.}
\label{tab:selector_effect}
\begin{tabular}{lrrrrr}
\toprule
Environment & $n$ & Archive HV & W/L/T & Final HV & W/L/T \\
\midrule
HalfCheetah & 15 & -0.15\% & 3/5/7 & -0.85\% & 4/4/7 \\
Hopper & 5 & +6.54\% & 2/1/2 & +59.56\% & 2/1/2 \\
Ant & 5 & -1.00\% & 1/3/1 & -3.33\% & 0/4/1 \\
\bottomrule
\end{tabular}
\end{table}

